\section{Notes on SDWA and its enforcement}
What changed around the ARP (1995) versus what changed earlier (1992–1993):

The most important timing issue is that the largest structural change in monitoring and reporting burden occurred in 1992–1993, not in 1995. This has two implications for your pre/post-1995 comparison:

1. Level shift in MR violation potential (1993, before the instrument):
Eighteen new SOC contaminants and ten new VOCs came under mandatory monitoring starting January 1, 1993. Each new contaminant creates an independent channel for MR violations (missed sample, late submission, missed confirmation sample, failure to certify TT for acrylamide/epichlorohydrin). This was a substantial expansion of administrative burden, but it is symmetric across all CWSs and pre-dates the ARP, so it enters the PWSID fixed effects as a within-unit level shift.

2. Compliance period transition coincides with the ARP (1996):
The second three-year compliance period runs 1996–1998, straddling the ARP. The monitoring schedule for that period depends on what was detected in the first period (1993–1995):

CWSs with no detections in 1993–1995 could reduce monitoring in 1996–1998 (SOC waivers, VOC annual or triennial samples). This decreases their MR exposure.
CWSs with detections in 1993–1995 were locked into quarterly monitoring in 1996–1998, increasing MR exposure.
If coal-mining areas were more likely to detect SOC or VOC contamination in 1993–1995 (plausible for certain pesticides), then the second compliance period creates a mechanically higher MR violation rate in mining areas in the post-1995 period, independent of the ARP. Conversely, if the ARP reduced production and thereby reduced contamination post-1995, it could have helped mining-area CWSs qualify for reduced monitoring in the third compliance period (1999–2001), reducing MR violations through a second-order channel.

3. What did NOT change around 1995:

Arsenic MCL was unchanged (remained 0.05 mg/l from the 1975 standard; new 0.01 mg/l rule was not promulgated until January 2001, effective January 2006)
Radionuclide standards were unchanged from their 1976 interim status
Total Coliform and SWTR requirements were already in effect since December 1990
No new IOC, SOC, or VOC MCLs took effect in or around 1995 from this rule


Surface Water Treatment Rule and Total Coliform Rule (54 FR 27486 and 54 FR 27544, both June 29, 1989): Required filtration and disinfection for surface water systems (TT approach for Giardia and viruses) and set a coliform MCL of no more than 5\% of monthly samples positive.

Thallium
Gross Alpha
Gross Beta 

The second three-year compliance period runs 1996–1998, straddling the ARP. The monitoring schedule for that period depends on what was detected in the first period (1993–1995): CWSs with no detections in 1993–1995 could reduce monitoring in 1996–1998 (SOC waivers, VOC annual or triennial samples). This decreases their MR exposure. CWSs with detections in 1993–1995 were locked into quarterly monitoring in 1996–1998, increasing MR exposure.

If coal-mining areas were more likely to detect SOC or VOC contamination in 1993–1995 (plausible for certain pesticides), then the second compliance period creates a mechanically higher MR violation rate in mining areas in the post-1995 period, independent of the ARP. Conversely, if the ARP reduced production and thereby reduced contamination post-1995, it could have helped mining-area CWSs qualify for reduced monitoring in the third compliance period (1999–2001), reducing MR violations through a second-order channel.

three amendments The list of contaminants expanded over time from 22 in 1976 to 91 in 2021. The first contaminants from 1976 were nitrates (10 mg/L as N), arsenic (0.05 mg/L), select inorganic chemicals (barium, cadmium, chromium, lead, mercury, selenium, silver), total coliforms, Synthetic Organic Chemicals (2,4-D, 2,4,5-TP (Silvex), Endrin, Lindane, Methoxychlor, Toxaphene), and radiological contaminants.

Between 1985 to 2005 the list of regulated contaminants and some of their MCL's changed. Arsenic. The MCL of 0.05 mg/L established in 1975 remained operative throughout the study period; the revised MCL of 0.010 mg/L under the 2001 Arsenic Rule (66 FR 6976, January 22, 2001) did not take effect until January 2006. Phase I Rule (52 FR 25690, 1987) established MCLs for eight VOCs, including benzene, carbon tetrachloride, trichloroethylene, vinyl chloride, and 1,2-dichloroethane. Surface Water Treatment Rule and Total Coliform Rule (54 FR 27486 and 54 FR 27544, both June 29, 1989): Required filtration and disinfection for surface water systems (TT approach for Giardia and viruses) and set a coliform MCL of no more than 5\% of monthly samples positive. Phase II Rule (56 FR 3526, January 30, 1991) and Phase V Rule (57 FR 31776, July 17, 1992): Established or revised MCLs for more than 50 SOCs (pesticides and PCBs) and additional inorganic chemicals. Between 1985 and 2005 arsenic and nitrate MCL's did not change and once established VOCs and SOCs MCL's did not change.

CWS self monitoring is the primary way the regulator determines if the SDWA requirements are being met. The frequency and location of water samples taken to fulfil the monitoring requirements changed several times between 1985 and 2005. Under the 1976 Interim Primary Drinking Water Standards (IPDWR) the number of microbial tests varied according to size of the population served, the water intake being underground or surface, and the number of previous positive tests. For inorganic chemicals, CWS using ground water only the testing required every three years testing and using surface water required yearly testing. A survey of Indiana CWS managers who responded to a 1980 survey on the impact of the SDWA saying that the SDWA's reporting requirements had the greatest impact on their systems, while MCL had the least (\cite{oleckno1982national}). Many contaminants required their own tests had their own methods which may be done on site or required a laboratory. The CWS had 40 days to report the results of a test to the state regulator and 48 hours to report a failure to meet the testing or monitoring standard. The records required submitted to the regulator required the name of the person doing the test, the date of the test, and the results of the test. Phase I Rule (52 FR 25690, July 8, 1987): all CWSs must sample quarterly. All systems that detect VOCs must continue testing quarterly. All ground water and surface water systems that are vulnerable to VOCs and do not detect VOCs in the first sample can reduce to testing every 5 years if they serve less than 500 connections, and 3 years if they serve more than 500 connections. Surface water systems that do not detect VOCs in the first year of quarterly sampling and are not at risk are only required to monitoring is only required at state discretion.

Phase II Rule (56 FR 3526, January 30, 1991) and Phase V Rule (57 FR 31776, July 17, 1992): Established or revised MCLs for more than 50 SOCs (pesticides and PCBs) and additional inorganic chemicals, and introduced the Standardized Monitoring Framework (SMF) organizing compliance into 9-year cycles of three 3-year periods. Except for microbial contaminants monitoring was 3 for surface and 1 for g 

The SDWA Monitoring and Enforcement Regime, 1985–2005
Monitoring and Reporting vs. Maximum Contaminant Level Violations


Regulatory Timeline and Contaminant-Specific Monitoring Requirements
1975 Interim Primary Standards and the original framework. EPA promulgated interim primary drinking water standards in December 1975 (40 FR 59566, December 24, 1975), establishing enforceable MCLs for nitrate (10 mg/L as N), arsenic (0.05 mg/L), select inorganic chemicals (barium, cadmium, chromium, lead, mercury, selenium, silver), organic chemicals, total coliforms, and radiological contaminants. These interim standards also referenced the 1976 Radionuclides Interim Rule (41 FR 28402, July 9, 1976), which set MCLs for combined radium-226+228 (5 pCi/L), gross alpha excluding radon and uranium (15 pCi/L), and beta/photon emitters (4 mrem/year). 



"The 1996 amendments strengthened enforcement authorities available to EPA and primacy states, streamlined the process for issuing federal administrative orders, increased administrative penalty amounts, made more sections of the act clearly subject to EPA enforcement, and required states (as a condition of primacy) to have administrative penalty authority. The amendments also provided that no enforcement action may be taken against a public water system that has a plan to consolidate with another system" (\cite{tiemann2014safe}).

"Citizen Suits Section 1449 provides for citizens’ civil actions. Citizen suits may be brought against any person or agency allegedly in violation of SDWA requirements or against the EPA Administrator for alleged failure to perform any action or duty that is not discretionary" (\cite{tiemann2014safe}).

Section 1414 also requires community water systems serving fewer than 10,000 people to mail to 
all customers an annual “consumer confidence report” on contaminants detected in their drinking 
water. Operators of public water systems serving more than 10,000 consumers are required to 
provide confidence reports biannually (\cite{tiemann2014safe})

SDWA 1996 sec. 105 provides technology recommendations to meet MCLs for systems of varying size. sec 1445 "In requiring a public water system to monitor under this subsection, the Administrator may take into consideration the system size and the contaminants likely to be found in the system's drinking water." Sec 1418 the state with primacy or the governor of a state without primacy can issue or apply for a reduced monitoring schedule on a contaminant by contaminant basis if the initial monitoring period finds no pollutants and the hydrogeology of the area supplying water does not contain sources of that pollutant. A public water system must show the state its source does not contain that pollutant. This does not apply to micorbial contaminants.

Relating to the 1996 SDWA amendments, "Section 1420 required states to establish capacity 
development programs, also based on EPA guidance. Congress specified that state programs 
include (1) legal authority to ensure that new systems have the technical, managerial, and 
financial capacity to meet SDWA requirements; and (2) a strategy to assist existing systems that 
are experiencing difficulties to come into compliance. EPA is required to withhold a portion of 
SRF grants from states that do not have capacity development strategies. The agency has not had 
to withhold funds under either of these programs." (\cite{tiemann2014safe})

Arsenic. The MCL of 0.05 mg/L established in 1975 remained operative throughout the study period; the revised MCL of 0.010 mg/L under the 2001 Arsenic Rule (66 FR 6976, January 22, 2001) did not take effect until January 2006. Under the SMF, groundwater systems sample once per 3-year compliance period, surface water systems annually, with waivers available to once per nine years for groundwater systems with sustained non-detects (40 CFR 141.23).

Inorganic chemicals (barium, chromium, etc.). Phase II and Phase V rules set MCLs for barium, beryllium, cadmium, chromium, cyanide, fluoride, mercury, nickel, selenium, and thallium. Monitoring follows the SMF: once per 3-year period for groundwater, annually for surface water, with waiver to 9-year cycles available (40 CFR 141.23).

Radionuclides. The 1977 interim standards governed monitoring through 1999; infrequent four-year monitoring cycles and weak enforcement resulted in sparse compliance data (CDC, History of Drinking Water Regulations). The final Radionuclides Rule (65 FR 76708, December 7, 2000) retained the 1977 MCLs for radium and gross alpha, added a uranium MCL of 30 µg/L effective December 2003, and introduced a tiered monitoring framework: systems with results below detection monitor once every nine years; those between 50\% and 100\% of the MCL every three years; those exceeding the MCL quarterly until four consecutive quarters return below the limit (40 CFR 141.26).

Total coliforms. Under the Total Coliform Rule (54 FR 27544, June 29, 1989, effective December 31, 1990), the required number of monthly samples scales with population served: from one sample per month for the smallest systems (25–1,000 persons) to 480 per month for the largest. Small groundwater systems with favorable sanitary survey history may be approved by the state for reduced monitoring of one sample per quarter. Any coliform-positive sample triggers repeat sampling of three to four additional samples within 24 hours (40 CFR 141.21).

Surface Water Treatment Rule. The SWTR operates on treatment technique rather than an MCL. Systems using surface water or groundwater under direct influence of surface water (GWUDI) must achieve 99.9\% (3-log) removal or inactivation of Giardia lamblia and 99.99\% (4-log) of viruses. Filtered systems must maintain turbidity \leq 0.5 NTU in 95\% of daily samples, never exceeding 5 NTU. Unfiltered systems must meet source water quality criteria and undergo state sanitary surveys at least every 18 months; the 1996 Amendments formalized surveys on a 3-year cycle for surface water systems. The Long Term 1 Enhanced SWTR (67 FR 1812, January 14, 2002) extended Cryptosporidium protections to small surface water systems. No federal Ground Water Rule was finalized until October 2006, outside the study window.

VOCs and SOCs. Under the Phase I rule (1987) and Phase II/V SMF (effective 1993), initial monitoring requires four consecutive quarterly samples at each entry point. With no detections, systems may reduce to annual monitoring and, for groundwater systems, ultimately to once per 3-year compliance period with a state-issued waiver (40 CFR 141.24). Systems with detections above 0.0005 mg/L must monitor quarterly until they demonstrate reliable compliance below the MCL. Surface water systems generally cannot receive waivers. SOC monitoring waivers are more commonly granted than VOC waivers because agricultural source profiles can be ruled out by vulnerability assessment.

What are the penalties for failing an MCL or MR under the SDWA (states may have penalties but the SDWA first uses technical assistance when an MCL violation occurs. When there is willful violation and no record keeping then there are civil actions to bring system in compliance and potentially \$5000 fine) Who enforces the penalties (the state EPA if they have primacy, the USEPA if the state fails, the courts).

A survey of indiana CWS managers in 1980 of the impact of the SDWA found that the Regulations' reporting requirements were found to have the greatest degree of impact on the systems, MCL the least (\cite{oleckno1982national}).

Monitoring for inorganic chemicals
has shown that 1500-3000 water supply
systems have levels above the NIPDWR
MCLs for certain of the contaminants.
These inorganics are primarily a problem
found in groundwaters, and removal of
inorganic chemicals can be difficult and
relatively expensive on a per capita
basis for small public water systems.
Problems continue primarily with compliance with the MCLs for arsenic,
barium, lead (from pipe or solder corrosion), radium, fluoride, and, to an increasing degree, nitrate. Small CWSs have struggled the most with implementing MCL and monitoring and reporting due to per capita costs and financial and human resource constraints.

The SDWA provides that if a system
will not be able to comply with an MCL after installation or use of the "best
technology, treatment techniques, or
other means, which the Administrator
finds to be generally available (GAT),"
taking costs into consideration, the system may apply for a variance (Section
1415 [a] [1] [A], 42 USC 300g-4 [a] [1] [A]).

burden on small systems or on those
states that conduct monitoring for certain systems (e.g., small systems) within
their boundaries. The states have reported that certain of these inorganic
compounds have not been detected at
significant levels in the drinking water
in many systems, and the probability of
future contamination is very slight.
Monitoring has shown that little change
in concentrations occurs over time for
certain contaminants, primarily inorganics in groundwater. In addition, some
contaminants, such as the six pesticides
named in the interim regulations, have
been found only rarely since compliance
monitoring requirements went into
effect.
To provide for more efficient use of
state and local resources, flexibility in
monitoring requirements will be a general principle in developing the revised
regulations. In addition, to assure detection and control of intermittent drinking
water contaminants (those that are not
homogeneously distributed), more specific monitoring requirements will be
designe

Because of the added costs of
treatment, many systems remain out of
compliance with various MCLs; some
systems remain unconvinced that the
net benefit of reducing contaminants is
worth the costs. This issue generally
relates directly to the availability and
strength of evidence contained in data
on potential health effect

Generally available technology (GAT)
would be defined for each regulated
contaminant, taking costs into consideration and possibly categorizing by
system characteristics, such as size or
water source. The states would evaluate
each case on a site-specific basis to
determine if the identified GAT was
appropriate and effective for that system.
In addition to central treatment alterna

1974 SDWA states get primacy partly when the state has adequate procedures to bring CWSs into compliance but allows discretion in the law about what that may be exactly (SDWA 1974 sec. 1413). When the state has primacy the actions taken against violators are escalitory and begin with the USEPA offering technical assitance and advice to bring the CWS into compliance, and allowing the state EPA to determine its own compliance plan. However, if the plan is unsatisfactory or does not lead to a compliance then the EPA steps in and the statute specificies USEPA can bring civil action against the CWS to protect public health and in the case where there was willful violation there is civil fine of no more than \$5000 on the CWS (SDWA 1974 sec. 1414). When a CWS is in violation of an MCL or reporting or has a variance for an MCL they must send notification to customers or face a \$5000 fine (SDWA 1974 sec 1414). 

CWSs have to maintain records, testing, and submit reports and tests that enable the administrator to determine if they are meeting the requirements of the act or face a fine no larger than \$5000 (SDWA 1974 sec. 1445).


Phase I Rule (52 FR 25690, July 8, 1987): Established MCLs for eight VOCs, including benzene, carbon tetrachloride, trichloroethylene, vinyl chloride, and 1,2-dichloroethane. Monitoring for VOCs initially quarterly in the first year to determine risk and prevelance, then state can reduce monitoring to quarterly to every 5 years based on what was detected and risk of system.

The 1986 SDWA Amendments (P.L. 99-339, June 19, 1986) sec 1445 for monitoring for unregulated contaminants, take into consideration the system size and the contaminants likely to be found in the system's drinking water. 

It also increased the penalty for violating a court order on a CWS to take action to protect public health from a max of \$5000 a day to \$25000 a day, also removed the requirement that the violation of the court order needs to be willful. The frequency of public notifications for exceeding an MCL or not monitoring depend on the health severity and frequency of the violation, the fine for not doing this increased from a max of \$5000 to \$25000. Congress directed EPA to promulgate NPDWRs for 83 specified contaminants by June 1989 and 25 more every three years thereafter (Congressional Research Service, R46652). This mandate drove three rapid rulemakings:

Phase I Rule (52 FR 25690, July 8, 1987): Established MCLs for eight VOCs, including benzene, carbon tetrachloride, trichloroethylene, vinyl chloride, and 1,2-dichloroethane. Monitoring for VOCs initially quarterly in the first year to determine risk and prevelance, then state can reduce monitoring to quarterly to every 5 years based on what was detected and risk of system. Systems with more than 10000 users have to begin monitoring by 1989 and those with less than 3300 begin by 1991. Each SDWA update requires states to bring their monitoring and testing and regulations up to the standard of the sdwa update to maintain primacy. CWSs are not required to use best available technology to meet MCL requirements but you need to use BAT to get a variance, also need to show a variance wont have high health risks.


Surface Water Treatment Rule and Total Coliform Rule (54 FR 27486 and 54 FR 27544, both June 29, 1989): Required filtration and disinfection for surface water systems (TT approach for Giardia and viruses) and set a coliform MCL of no more than 5\% of monthly samples positive.

Phase II Rule (56 FR 3526, January 30, 1991) and Phase V Rule (57 FR 31776, July 17, 1992): Established or revised MCLs for more than 50 SOCs (pesticides and PCBs) and additional inorganic chemicals, and introduced the Standardized Monitoring Framework (SMF) organizing compliance into 9-year cycles of three 3-year periods.
Under SMF base testing for inorganics was every 3 years for ground water and once a year for surface water - the first one had to be done by end of 1995. The frequency could change based on system size and state determination of risks. VOCs must be tested quarterly at first and then for every year for surface water and every three years for ground water if state allows them an exemption and they did not detect VOCs - detecting VOCs became a condition for getting reduced testing. Under the SMF, surface water-supplied CWSs monitor quarterly (four samples per year) at each entry point; groundwater systems monitor annually. States may grant reduced monitoring to annual for surface water systems with four consecutive quarterly results below 50\% of the MCL, and groundwater systems sustaining prolonged compliance may be reduced to one sample per nine-year cycle (40 CFR 141.23).

Phase V NPDW 1992 VOCs monitor quarterly then annually unless a groundwater system is granted an exemption by the state to test every 3 years. If they detect an MCL or some level of VOC then they have to monitor quarterly. SOCs quarterly and then four consecutive quarter samples every 3 years. Or two consecutive quarter samples every three years if no SOC detected and state determines low risk. If any SOC detected then back to quarterly sampling.


The 1996 SDWA Amendments (P.L. 104-182, August 6, 1996) replaced the mandatory 25-contaminants-every-three-years schedule with a risk-based, cost-benefit framework (Centers for Disease Control and Prevention, History of Drinking Water Regulations). Key changes relevant to monitoring and enforcement: the per-violation civil penalty ceiling rose from \$5,000/day to \$25,000/day; streamlined administrative order authority eliminated prior public-hearing requirements before EPA could issue orders; states were required to have administrative penalty authority as a condition of maintaining primacy; and the Standardized Monitoring Framework was supplemented by Alternative Monitoring Guidelines allowing states to modify groundwater monitoring frequencies for non-vulnerable systems (42 USC §300g-3; Congressional Research Service, R46652).

Enforcement Authority and Actions
SDWA enforcement responsibility is primarily delegated to state primacy agencies; as of 2005, 55 of 57 states and territories held primacy for public water system supervision (Congressional Research Service, R46652). To obtain primacy, a state must adopt regulations at least as stringent as NPDWRs, maintain inspection and penalty authority, conduct sanitary surveys, and report violations to EPA.

The formal enforcement sequence begins with detection — either through system self-reporting or state inspection — followed by informal compliance assistance (letters, notices of violation) and, if violations persist, formal administrative actions. Under 42 USC §300g-3, when EPA determines that a state has not commenced enforcement within 30 days of notification of a violation, EPA must issue a federal administrative compliance order (ACO) or commence a civil action. In non-primacy states, EPA acts directly without the 30-day waiting period.

Civil penalties before 1996 were capped at \$5,000 per violation per day; the 1996 Amendments raised this to \$25,000 per day (adjusted for inflation in subsequent years). Administrative penalties below \$5,000 require notice and public hearing; those between \$5,000 and \$25,000 require a formal on-the-record hearing; penalties exceeding \$25,000 proceed to federal district court via DOJ referral. EPA referred on average approximately 420 cases per year to DOJ for civil or criminal judicial enforcement between 1998 and 2011; total civil judicial penalties averaged approximately \$93 million per year (Shimshack 2014, p. 344). Ordinary SDWA noncompliance does not trigger criminal liability; criminal provisions are reserved for tampering with a water system (up to 20 years imprisonment, 42 USC §300i-1) and willful underground injection control violations.

Empirical evidence indicates that formal penalties generate substantial general deterrence: Shimshack and Ward (2005) find that a single enforcement fine reduces the statewide violation rate by roughly two-thirds in the following year, primarily through spillover effects to uninspected systems. Administrative orders without financial penalties show no detectable compliance effect, suggesting that the monetary component is the operative deterrent (Shimshack and Ward 2005, 2008).

States and EPA also have several flexibility mechanisms that reduce enforcement burden: variances (available to systems whose raw water quality prevents MCL compliance despite best technology), exemptions (up to nine years for small systems serving 3,300 or fewer), and — after the 1996 Amendments — compliance schedules tied to consolidation or ownership transfer plans (42 USC §300g-3(h); Congressional Research Service, R46652).

Key references:

Congressional Research Service. The Safe Drinking Water Act (SDWA): A Summary of the Act and Its Major Requirements, R46652.
Centers for Disease Control and Prevention. History of Drinking Water Regulations. (CDC, Healthy Water).
[SDWA history.pdf — CRS/sdwa history document, project lit folder]
Shimshack, J.P. (2014). "The Economics of Environmental Monitoring and Enforcement." Annual Review of Resource Economics 6: 339–360.
Shimshack, J.P. and Ward, M.B. (2005). "Regulator Reputation, Enforcement, and Environmental Compliance." Journal of Environmental Economics and Management 50(3): 519–540.
Shimshack, J.P. and Ward, M.B. (2008). "Enforcement and Over-Compliance." Journal of Environmental Economics and Management 55(1): 90–105.
Helland, E. (1998). "The Enforcement of Pollution Control Laws: Inspections, Violations, and Self-Reporting." Review of Economics and Statistics 80(1): 141–153.
40 CFR Part 141, Subpart C (monitoring requirements); 42 USC §300g-3 (enforcement).
Federal Register citations for individual rules are given inline above.
A few notes on using this in your paper: (1) The arsenic MCL was still 0.05 mg/L throughout 1985–2005, so arsenic MCL violations in your data reflect the old standard. (2) Radionuclide monitoring was genuinely sparse and state-directed pre-2000, which may explain thinner violation data for radionuclides in your early sample years. (3) VOC and SOC monitoring did not begin in earnest until Phase I (1987) and SMF rollout (1993), so those outcomes are effectively unidentified before ~1993. Worth noting when reporting sample-period results for those contaminants.